\section{Ablação e diagnóstico das camadas intermediárias}

A análise complementar é retrospectiva e utiliza o mesmo motor, custos de 10 bps e rebalanceamento mensal. As cardinalidades 5, 10 e 20 são ilustrativas. Não houve nova seleção de hiperparâmetros.

\begin{table}[H]\centering\caption{Ablação com seleção própria em 2025.}

\begin{tabular}{lrrrr}\hline Variante & n & Retorno (\%) & Sharpe & MDD (\%) \\ \hline

Sem views & 5 & 59.24 & 1.60 & -10.73 \\

Sem views & 10 & 42.63 & 1.22 & -10.54 \\

Sem views & 20 & 28.21 & 0.74 & -8.46 \\

XGB com gate & 5 & 40.26 & 1.02 & -16.30 \\

XGB com gate & 10 & 33.78 & 0.87 & -13.82 \\

XGB com gate & 20 & 20.30 & 0.35 & -13.67 \\

XGB livre & 5 & -0.89 & -0.70 & -16.29 \\

XGB livre & 10 & -0.70 & -0.69 & -16.12 \\

XGB livre & 20 & -0.41 & -0.68 & -15.82 \\

Textual & 5 & 51.48 & 1.41 & -11.94 \\

Textual & 10 & 44.08 & 1.26 & -9.84 \\

Textual & 20 & 26.53 & 0.63 & -9.89 \\

Híbrido & 5 & 58.40 & 1.60 & -10.73 \\

Híbrido & 10 & 40.14 & 1.13 & -10.20 \\

Híbrido & 20 & 28.21 & 0.73 & -8.46 \\

\hline\end{tabular}\par\smallskip{\footnotesize Fonte: elaboração própria.}\end{table}

A variante XGB com gate preserva a confiança híbrida e remove somente a parcela textual de Q. Sua comparação com o híbrido avalia o efeito incremental dessa parcela sob a mesma regra de cobertura. XGB livre remove também a dependência de notícias e usa confiança quantitativa. A variante textual utiliza confiança de notícias, mas ainda herda escala e horizonte quantitativos; portanto, não representa um sistema textual inteiramente independente.

Nos conjuntos comuns, as comparações são pareadas por subconjunto; os dados completos estão nos artefatos de auditoria. Com vinte ativos, elimina-se a diferença de elegibilidade entre as estratégias. O resultado descreve o efeito das intervenções neste motor e neste período, sem demonstrar generalização futura.

O resultado mais relevante é a proximidade entre o híbrido e o prior sem views com vinte ativos: 28,2085\% e 28,2053\% de retorno líquido, respectivamente, diferença de apenas 0,0032 ponto percentual. Com seleção própria de cinco e dez ativos, o prior retorna 59,24\% e 42,63\%, frente a 58,40\% e 40,14\% do híbrido. Não se observa ganho consistente de retorno ao adicionar as views ao prior nessas trajetórias.

Ao manter o gate e a confiança híbrida, mas retirar o sentimento de Q, os retornos passam a 40,26\%, 33,78\% e 20,30\% para cinco, dez e vinte ativos. A parcela textual melhora o resultado frente a esse controle, mas isso não demonstra superioridade frente à ausência de views. O XGB sem gate apresenta retornos negativos nas três cardinalidades; essa variante altera também a confiança e deve ser interpretada como intervenção conjunta.

A cobertura de seis pares habilitados permanece uma medida de elegibilidade, inclusive nas tabelas de auditoria das variantes sem gate. O afastamento entre posterior e prior e a distância entre pesos-alvo foram exportados para auditoria. A meia distância L1 entre pesos-alvo não equivale ao turnover realizado, pois ignora a deriva dos pesos entre negociações.

Esses resultados reforçam que a vantagem frente a Markowitz não pode ser atribuída automaticamente ao aprendizado de máquina ou ao sentimento. O prior, a estrutura de covariância e as restrições respondem por parte substancial do comportamento observado. A interpretação continua limitada por escolhas retrospectivas e por apenas um ano de avaliação financeira.
